\tocsection{\textbf{10/27/26 Lecture 18:} Beyond James-Stein and Minimax Estimation}

\textbf{Recall:} For $X \sim N_d(\theta, I_d)$, for the shrinkage estimator $(1 - \zeta)X$, we found
\[ \hat{\zeta}_{\tt{UMVU}} = \frac{d - 2}{\norm{X}^2} \rightsquigarrow \delta_{\tt{JS}} = \left(1 - \frac{d - 2}{\norm{X}^2}\right)X \]
for $\zeta = \frac{1}{1 + \tau^2}$, but even without a prior on $\tau$, $\delta_{\tt{JS}}$ still is better than $X$. We also had
\[ \tt{MSE}_{\theta}(\delta_{\tt{JS}}) = d - \E_{\theta}\left[\frac{(d - 2)^2}{\norm{X}^2}\right] \]
and $\tt{MSE}_0(\delta_{\tt{JS}}) = 2$. Plotting,
\fig{0.42}{Lecture18Fig1}{}

Now, note that $\frac{d - 2}{\norm{X}^2}$ can be $\geq 1$, so an even better estimator is
\[ \delta_{\tt{JS}^{+}}(X) = \left(1 - \frac{d - 2}{\norm{X}^2}\right)_{+}X \]
Another thing you can do is consider
\[ \delta_{\tt{JS}, 2}(X) = \left(1 - \frac{d - 3}{\norm{X - \overline{X}1_d}^2}\right)(X - \overline{X}1_d) + \overline{X}1_d \]
for $d \geq 4$.\\

For $d \geq 4$, $\delta_{\tt{JS}, 2}(X)$ strictly dominates $\delta_0$.\\

\textbf{Idea of derivation:} Change of basis so that $\overline{X}1_d$ is a basis vector, and apply JS on the $(d - 1)$-dimensional subspace orthogonal to it, using the projection onto the subspace.

\subsection*{Minimax Risk}
\sectionspace

The idea is to find
\[ \underset{\delta}{\tt{minimize}}\left(\sup_{\theta} R(\theta, \delta)\right) \]

\begin{definition}[\textbf{Definition 18.1:} Minimax Risk]
    The \emph{minimax risk} is $r^{*} = \inf_{\delta}\sup_{\theta} R(\theta, \delta)$. An estimator $\delta^{*}$ is \emph{minimax} if $\sup_{\theta} R(\theta, \delta^{*}) = r^{*}$.
\end{definition}

In a sense, this is a very pessimistic estimator, but it has some uses. For example, in game theory, if someone was picking a purposely difficult-to-estimate $\theta$, $\delta^{*}$ is good. It can also be used to bound how hard estimation of a problem is.\\

This is also connected to Bayes, as $\delta^{*}$ is okay for all priors. That is, if there is some prior that you don't know, then $\delta^{*}$ is good.\\

\textbf{Note:} For any proper prior $\Lambda$, the Bayes risk is
\begin{align*}
    r_{\Lambda} &= \inf_{\delta}\int R(\theta, \delta) d\Lambda(\theta)\\
    &\leq \inf_{\delta}\sup_{\theta} R(\theta, \delta)\\
    &= r^{*}
\end{align*}
Moreover, if $\delta_{\Lambda}$ is Bayes, then $r_{\Lambda} = \int R(\theta, \delta_{\Lambda}) d\Lambda(\theta)$, so the Bayes risk of any prior lower bounds $r^{*}$.

\begin{definition}[\textbf{Definition 18.2:} Least Favorable Priors]
    A \emph{least favorable prior} $\Lambda^{*}$ gives the best lower bound
    \[ r_{\Lambda^{*}} = \sup_{\Lambda} r_{\Lambda} \]
    so for any $\Lambda$ and any $\delta$,
    \[ r_{\Lambda} \leq r_{\Lambda^{*}} \leq r^{*} \leq \sup_{\theta} R(\theta, \delta) \]
\end{definition}

\begin{theorem}[\textbf{Theorem 18.3}: Correspondence between Bayes and Minimax]
    Let $\Lambda$ be a prior with Bayes estimator $\delta_{\Lambda}$. If $r_{\Lambda} = \sup_{\theta} R(\theta; \delta_{\Lambda})$, then $\delta_{\Lambda}$ is minimax, $\Lambda$ is LF, and $r^{*} = r_{\Lambda}$.\\

    Moreover, if $\delta_{\Lambda}$ is the unique Bayes estimator up to $\overset{\tt{a.s.}}{=}$, it is also the unique minimax estimator.
\end{theorem}
\begin{proof}{Proof:}
    The first statement is obvious from above:
    \[ \sup_{\theta} R(\theta, \delta_{\Lambda}) = r_{\Lambda} \leq r^{*} \leq \sup_{\theta} R(\theta, \delta_{\Lambda}) \]
    For any other $\delta$,
    \begin{align*}
        \sup_{\theta} R(\theta, \delta) \geq \int R(\theta, \delta) d\Lambda(\theta) &> \int R(\theta, \delta_{\Lambda}) d\Lambda(\theta)\\
        &= r_{\Lambda} = r^{*}
    \end{align*}
    where $>$ comes from uniqueness of $\delta_{\Lambda}$.

    $\qed$
\end{proof}

\textbf{Note:} The condition $r_{\Lambda} = \sup_{\theta} R(\theta; \delta_{\Lambda})$ implies we have almost everywhere constant risk on the support of $\Lambda$.\\

\textbf{Checkable Conditions:}
\begin{enumerate}
    \item $\delta_{\Lambda}$ is a Bayes estimator with constant risk.
    \item $\delta_{\Lambda}$ is Bayes and its risk attains its max on the support of $\Lambda$.
\end{enumerate}

\textbf{Example:} $X \sim N(\theta, 1)$, $\abs{\theta} \geq 1$. We want to estimate $g(\theta) = \tt{sgn}(\theta)$.\\

Intuitively, $\Lambda = \tt{Unif}(\{-1, 1\})$ is a pretty unfavorable prior, so let's try to find a Bayes estimator for this.
\[ \bb{P}_{\Lambda}(\theta = 1 \given X = x) = \frac{\phi(x - 1)}{\phi(x - 1) + \phi(x + 1)} \quand \bb{P}_{\Lambda}(\theta = -1 \given X = x) = \frac{\phi(x + 1)}{\phi(x - 1) + \phi(x + 1)} \]
so
\[ \E_{\Lambda}(g(\theta) \given X = x) = \frac{\phi(x - 1) - \phi(x + 1)}{\phi(x - 1) + \phi(x + 1)} = \frac{e^{x} - e^{-x}}{e^{x} + e^{-x}} = \tanh(x) \]
Then, we see
\fig{0.42}{Lecture18Fig2}{}

Then, as the risk of $\delta_{\Lambda}$ is constant on the support of $\Lambda$ (and maximized there), we see $\delta_{\Lambda}$ is minimax.\\

\textbf{Example:} $X \sim \Binom(n, \theta)$, estimate $\theta$ using MSE.\\

Consider $\theta \sim \tt{Beta}(\alpha, \beta)$ and $\alpha = \beta$, so $\delta_{\alpha}(X) = \frac{X + \alpha}{n + 2\alpha}$. Now, consider
\[ \E_{\theta}\delta_{\alpha}(X) = \frac{n\theta + \alpha}{n + 2\alpha} \implies \tt{Bias}(\theta, \delta_{\alpha}) = \frac{\alpha(1 - 2\theta)}{n + 2\alpha} \]
\[ \Var_{\theta}(\delta_{\alpha}(X)) = \frac{n\theta(1 - \theta)}{(n + 2\alpha)^2} \]
Then

\begin{align*}
    \tt{MSE}(\delta_{\alpha}) &= \tt{Bias}^2 + \Var\\
    &= \frac{1}{(n + 2\alpha)^2}\left(\alpha^2(1 - 2\theta)^2 + n\theta(1 - \theta)\right)\\
    &= (n + 2\alpha)^{-2}\left(\alpha^2 + (n - 4\alpha^2)\theta(1 - \theta)\right)
\end{align*}

So we can take $\alpha = \frac{\sqrt{n}}{2}$, which makes this constant in $\theta$. Thus, $\delta_{\alpha}(X) = \frac{X + \sqrt{n}/2}{n + \sqrt{n}}$ is minimax. However, note this isn't very good by this graph:
\fig{0.45}{Lecture18Fig3}{}
